\section{Problem}

In computer graphics, realistic surface rendering is essential for high visual quality. A common technique is texture mapping, where image-based textures are applied to geometric surfaces to simulate materials such as wood, stone, or fabric. However, this approach has notable limitations: fixed image textures lack flexibility, scalability, and parameterization, making them less suitable for dynamic or procedurally generated environments.

Algorithmic texture synthesis offers an alternative by generating textures through computational methods rather than relying on pre-existing images. These include procedural approaches based on mathematical functions, rule-based methods that define pattern formation, and optimization-based techniques that estimate parameters to approximate a given target texture.

A central challenge is that algorithmically generated textures cannot exactly reproduce real-world samples due to structural differences. Consequently, texture approximation is formulated as an optimization problem, where model parameters are adjusted to maximize similarity to a target texture according to a defined metric.

This thesis therefore investigates algorithmic texture synthesis methods and their ability to approximate target textures, with a focus on procedural, rule-based, and optimization-based approaches, as well as their implementation, parameterization, and evaluation in computer graphics applications.

\section{Research Questions}

Based on the problem described above, this thesis aims to investigate several research questions related to the generation and approximation of textures through algorithmic methods.

The main research question can be formulated as follows:

\textit{How can algorithmic texture synthesis methods be applied in procedural real-time environments to improve variation, controllability, and scalability of generated textures?}

To address this main question, several sub-questions can be defined:

\begin{enumerate}
    \item What existing approaches and systems for algorithmic texture synthesis are described in current research and industry applications?
    \item How can procedural and rule-based texture synthesis methods be implemented within a procedural real-time environment?
    \item To what extent can procedural parameters be optimized in order to approximate target textures?
    \item How do different texture synthesis approaches compare with respect to visual quality, runtime performance, reproducibility, and parameter controllability?
\end{enumerate}

These research questions guide the investigation of algorithmic texture synthesis approaches and provide a structured framework for implementing, evaluating, and comparing the developed methods within the proposed research prototype.

\section{Vision}

The goal of this work is to advance the understanding of algorithmic texture generation and its practical applicability in computer graphics. By exploring procedural, rule-based, and optimization-based methods, the thesis demonstrates how textures can be generated in a flexible, parameter-driven manner instead of relying on static image resources.

Such approaches enable more scalable and adaptable workflows in areas including real-time graphics, game development, simulation, and procedural world generation. Algorithmic methods support dynamic texture creation, parameter-based modification, and context-dependent adaptation without requiring large texture libraries.

To support this vision, the thesis proposes a prototype system for implementing and evaluating parameterizable texture models. The results aim to provide insights into the strengths and limitations of different synthesis approaches and assess their suitability for typical computer graphics applications.

Ultimately, this work contributes to ongoing research in procedural content generation, where algorithmic techniques are increasingly used to create complex and visually rich environments.